\section{Discussion}

Our central finding is a sizable firm-size wage premium that survives controls
for gender, age, age squared, education, formality, and sector-year conditions. While the
magnitude of the premium is broadly consistent with the existing Latin American
evidence on employer-side wage differences \citep{Schaffner1998,
AlvarezBenguriaEngbomMoser2018,GerardLagosSeverniniCard2021}, the persistence
of the premium after controlling for formality is the more substantively
informative feature of the results. This section discusses what such a premium
can and cannot tell us about the underlying mechanisms.

A long-standing concern in the developing-country literature is that the
apparent firm-size premium is, at least in part, a relabeling of the
formal-informal wage gap. In dual labor-market models in the tradition of
\citet{HarrisTodaro1970} and \citet{Fields1975}, large firms are
disproportionately formal, and formality itself carries a wage premium through
some combination of higher productivity, mandated benefits, payroll taxation,
and selection. From this perspective, much of what is conventionally labeled a
``firm-size'' premium may simply reflect the boundary between the formal and
informal sectors rather than any return to firm scale per se, since formal firms
also tend to be larger.

By showing that the premium remains after conditioning on formality, our
estimates indicate that this dual-market channel is not the whole story. Even
among workers with the same formality status, larger employers pay more. In this
respect, our findings reinforce \citet{Schaffner1998} and complement the
matched employer-employee evidence from Brazil
\citep{AlvarezBenguriaEngbomMoser2018,GerardLagosSeverniniCard2021,
EngbomMoser2022}, which points toward employer-side pay determination as a
first-order feature of Latin American labor markets. Related firm-to-firm
evidence from Costa Rica shows that productive linkages with multinationals can
shift domestic firm performance and employment
\citep{AlfaroUrenaManeliciVasquez2022}.

What our controls cannot do is distinguish among other leading explanations for
the premium. As \citet{BrownMedoff1989} and \citet{OiIdson1999} emphasized for
the United States, a premium that survives observable controls is consistent
with several non-exclusive mechanisms. The first is unobserved worker quality:
gender, schooling, and sector capture only part of productive heterogeneity,
leaving non-cognitive skills, school quality, and unmeasured ability bundled
into the firm-size coefficient. At the same time, the available developing-
country evidence cautions against treating sorting as the whole explanation:
\citet{SoderbomTealWambugu2005} show that the firm-size effect remains
substantial in Kenya and Ghana even after controlling for individual fixed
effects.

A second possibility is rent sharing. Larger firms tend to be more productive,
more capital-intensive, and more exposed to product-market power, and they may
share part of those rents with workers through bargaining, fair-wage norms, or
hold-up considerations. The matched-data literature beginning with
\citet{AbowdKramarzMargolis1999} and \citet{CardHeiningKline2013}, and extended
to the United States by \citet{Song2019}, has shown that firm-specific pay
premia account for a quantitatively important share of wage variation in
advanced economies; the Brazilian evidence cited above extends this conclusion
to Latin America. A third possibility is efficiency-wage or monitoring-cost
reasoning, in which larger firms face greater difficulty observing individual
effort and pay above-market wages to elicit it. Because our identification rests
on observables, we cannot adjudicate among these mechanisms, and we do not
interpret the surviving coefficient as a pure rent-sharing parameter.

%One reading we explicitly do not draw is a direct monopsony test. The recent literature on firm wage setting emphasizes that employers may have wage-setting power, but monopsony concerns markdowns relative to marginal products and firm-specific labor-supply elasticities, neither of which is observed in the GEIH \citep{Card2022}. A positive, robust firm-size premium is therefore compatible with monopsonistic labor markets, but it is not evidence for or against monopsony by itself. In this paper, the premium is more directly informative about productivity-driven pay differences, rent sharing, worker sorting, and segmentation, consistent with the firm-centric view of wage determination that has emerged in the recent literature.

Several limitations follow from this discussion. Because our identification rests on observables, we cannot separate unobserved worker quality from the effect of firm scale. Doing so would require either matched employer-employee data with worker mobility, in the spirit of \citet{AbowdKramarzMargolis1999} and \citet{CardHeiningKline2013}, or a source of plausibly exogenous variation in firm-size exposure. Similarly, while controlling for formality rules out a purely mechanical formal-informal explanation, it does not rule out firm characteristics correlated with both size and pay-setting capacity, such as capital intensity, exporter status, or foreign ownership. These extensions are natural next steps. In the meantime, our results contribute to a growing body of evidence that the firm-size wage premium in Latin America is significant, robust to standard observable controls, and not reducible to the formal-informal margin.

The evidence does not imply that policy should simply push workers into large firms or favor large firms as such. The paper does not observe firm productivity or worker-firm matches, so it cannot prove that the wage premium is a causal return to scale. The policy relevance of the results is more precise: a large share of Colombian employment remains in solo work and micro-firms, the middle of the firm-size distribution is thin, and larger firms pay higher wages even after conditioning on observed worker characteristics, formality, and sector-year conditions. This combination is a warning sign of labor-market segmentation and possible allocative frictions.

The policy implication is that policies should avoid preserving smallness as an objective in itself. Support for small firms is most valuable when it helps productive firms grow, formalize, adopt better technologies, and connect to larger markets. By contrast, policies that subsidize informality or create thresholds that make growth costly may reinforce the same employment structure documented in the paper: many workers in very small units, few workers in the middle, and a large wage gap between low-scale and large-scale employment.